\subsection{Do classifiers gain from bone-suppressed inputs? (H6--H10)}
\label{sec:cls_results}
All 175 runs finished; Table~\ref{tab:cls_macro} gives the test metrics of the seed ensembles.

\begin{table}[htbp]
\centering\small
\caption{Test macro-AUROC and macro-AUPRC of the seed ensembles (five seeds; BDCXR-3257: out-of-fold over five folds). Descriptive; the
tests are in Tables~\ref{tab:cls_primary} and~\ref{tab:cls_groups}. orig: radiograph; soft: soft-tissue image; early: three images as
channels; hybrid: three streams; ctrl: radiograph in three streams.}
\label{tab:cls_macro}
\fittable{tables/v3_macro.tex}
\end{table}

\subsubsection{Overall performance (H6, H7)}
The hybrid did not improve on the radiograph alone on any dataset (Table~\ref{tab:cls_primary}, Fig.~\ref{fig:cls_forest}). On NIH the
macro-AUROC difference was $+0.0001$ $[-0.0022,0.0024]$ (Holm $p=0.967$), so a gain larger than 0.0024 is not compatible with the data; on
VinBigData it was $+0.0024$ $[-0.0004,0.0053]$ (Holm $p=0.28$) and on BDCXR-3257 $-0.0042$ $[-0.0108,0.0021]$ (Holm $p=0.39$). The
AUPRC differences agreed in direction and none excluded zero. The hybrid did not differ from the capacity control on any dataset (Holm
$p=1$). Against channel stacking it was nominally worse on NIH ($-0.0031$ $[-0.0055,-0.0007]$, Holm $p=0.066$; AUPRC $-0.0055$
$[-0.0100,-0.0012]$) and no different elsewhere.

\begin{table}[htbp]
\centering\footnotesize
\caption{H6 (primary) and H7: differences of the seed ensembles with 95\,\% cluster-bootstrap intervals (\num{10000} resamples). Holm: H6
over the three datasets, H7 over its six tests.}
\label{tab:cls_primary}
\fittable{tables/v3_primary.tex}
\end{table}

\begin{figure}[htbp]
\centering
\safefig{figures/fig_hypothesis_forest_densenet121.pdf}
\caption{Pre-specified contrasts with 95\,\% cluster-bootstrap intervals. (a) H6 and H7: macro-AUROC differences per dataset. (b) H8:
differences in the mean AUROC of each finding group's classes. Filled markers: Holm $p<0.05$; right: Holm-adjusted $p$.}
\label{fig:cls_forest}
\end{figure}

\subsubsection{Finding groups and classes (H8)}
The effect of the inputs depended on the finding group (Table~\ref{tab:cls_groups}, Fig.~\ref{fig:cls_forest}b). On NIH the soft-tissue
image raised the mean AUROC of the focal findings by $+0.0089$ $[0.0040,0.0140]$ (Holm $p=0.0024$) and the hybrid by $+0.0063$
$[0.0005,0.0117]$ (Holm $p=0.101$). Diffuse findings lost with both inputs (hybrid $-0.0052$ $[-0.0077,-0.0027]$, Holm $p=0.0024$; soft
$-0.0032$ $[-0.0057,-0.0008]$, Holm $p=0.039$), and boundary findings did not change. On VinBigData the focal estimates had the same direction
and size (hybrid $+0.0085$ $[-0.0023,0.0196]$; soft $+0.0083$ $[-0.0044,0.0202]$) but wide intervals, the diffuse estimates were close to
zero (hybrid $+0.0011$, soft $+0.0003$), and no group difference survived the correction.

\begin{table}[htbp]
\centering\footnotesize
\caption{H8: difference in the mean per-class AUROC of each finding group, with 95\,\% cluster-bootstrap intervals; Holm within dataset over
groups and contrasts.}
\label{tab:cls_groups}
\fittable{tables/v3_groups.tex}
\end{table}

Per class (Table~\ref{tab:cls_per_class}, Fig.~\ref{fig:cls_per_class}), two NIH classes differed between the hybrid and the
radiograph after correction: nodule gained ($+0.0129$, Holm $p=0.00121$) and infiltration lost ($-0.0160$, Holm $p<0.001$; also $-0.0128$
against channel stacking). Mass did not change ($-0.0004$). No class differed on VinBigData or BDCXR-3257.

\begin{table}[htbp]
\centering\footnotesize
\caption{H8 per class: test AUROC of the radiograph and hybrid seed ensembles and the paired DeLong contrasts of the hybrid against the
radiograph and against channel stacking; Holm within dataset over classes and both contrasts.}
\label{tab:cls_per_class}
\fittable{tables/v3_per_class.tex}
\end{table}

\begin{figure}[htbp]
\centering
\safefig{figures/fig_auroc_per_class_densenet121.pdf}
\caption{NIH ChestX-ray14: test AUROC per class and per finding group for each configuration, and each configuration minus the radiograph,
with 95\,\% patient-cluster bootstrap bands from 1,000 resamples (display only; the bootstrap tests of Tables~\ref{tab:cls_primary}
and~\ref{tab:cls_groups} use \num{10000}).}
\label{fig:cls_per_class}
\end{figure}

\subsubsection{Where the streams point (H10)}
Across 984 targets (164 focal) on 880 test images, Grad-CAM++ maps of the radiograph stream hit focal targets less often than other targets
(0.213 against 0.374), whereas soft-tissue stream maps hit both about equally (0.385 against 0.401) (Table~\ref{tab:cls_pointing},
Fig.~\ref{fig:cls_pointing}). The soft-minus-radiograph difference was therefore larger for focal targets ($+0.172$) than for the others
($+0.027$): interaction $+0.145$ $[0.092,0.199]$, $p=0.0002$. Integrated Gradients ranked the two streams the same way (focal: 0.580
soft-tissue against 0.505 radiograph; other: 0.621 against 0.645). Bone-stream maps rarely hit any target (0.049 and 0.167).
Fig.~\ref{fig:cls_maps} shows examples.

\begin{table}[htbp]
\centering\footnotesize
\caption{H10: pointing-game hit rates on the NIH expert boxes, means over targets of the per-target means over the five hybrid models;
tolerance 15\,px at 256\,px. Last row: soft-tissue minus radiograph stream (Grad-CAM++) and the interaction with its 95\,\% image-cluster
bootstrap interval.}
\label{tab:cls_pointing}
\fittable{tables/v3_pointing.tex}
\end{table}

\begin{figure}[htbp]
\centering
\safefig{figures/fig_pointing_game_densenet121.pdf}
\caption{Pointing game per NIH box class and per group. (a) Hit rate of each map. (b) Decoder routing: share of each class query's
attention on each stream. (c) Share of the Integrated-Gradients attribution per input. Brackets: number of targets.}
\label{fig:cls_pointing}
\end{figure}

\begin{figure}[htbp]
\centering
\safefig{figures/fig_xai_comparison_densenet121.pdf}
\caption{Attribution maps of one hybrid model (seed 0) on NIH ChestX-ray14 test images: the three inputs, Grad-CAM++ per stream and the
decoder attention, with the expert boxes. Rows: targets (row label: the box class).}
\label{fig:cls_maps}
\end{figure}

\subsubsection{Uncertainty}
Conformal risk control held the NIH false-negative rate at 0.060--0.069 on the test split for nominal 0.10, but the label sets held 6.33--6.78
of 14 labels on average (Supplementary Table~S10). On BDCXR-3257 the Mondrian sets covered 0.897--0.900 of the images at nominal 0.90
with 1.76--1.81 of 3 classes per set. In selective classification (Supplementary Table~S11 and Fig.~S13), adding the suppressor's spread to
the confidence score raised the AURC in every configuration: NIH 0.0352--0.0358 with the entropy alone against 0.0372--0.0375 with the spread,
both on the \num{25580} test images with a lung-field spread (16 of \num{25596} have an empty lung mask), and BDCXR-3257 0.1754--0.1938
against 0.2138--0.2260.

The supplementary material shows every single run next to the seed ensembles (Supplementary Fig.~S1), the conformal sets per
class (Supplementary Fig.~S2) and the Integrated-Gradients maps (Supplementary Fig.~S3).
